\section{Related Work}
\label{sec:related}

\para{Entity tracking benchmarks.}
\citet{kim2023entity} introduced the \boxes benchmark, a synthetic task requiring models to track objects across 7~boxes after a series of move operations.
They showed that earlier claims of entity tracking ability were inflated by trivial baselines, and that code-pretrained models substantially outperformed standard language models.
\citet{kim2024code} extended this analysis, confirming that code pretraining causally improves entity tracking: Code Llama 70B-Instruct achieved 64.9\% accuracy on non-trivial \boxes examples, the best result among open-weight models at the time.
Both studies fix the number of entities at~7 and vary the number of operations, leaving the question of how performance scales with entity \emph{count} unanswered.
The \ruler benchmark~\citep{hsieh2024ruler} includes a variable-tracking task analogous to entity tracking but uses chain-style assignments rather than narrative framing and does not systematically explore high entity counts.

\para{Narrative understanding and character tracking.}
\citet{hamilton2025tldm} evaluated frontier models on full novels, finding that all models lose stable narrative understanding beyond 64K tokens, with character-location tracking (``storyworld description'') degrading fastest among their three tasks.
However, because real novels confound character count with text length, discourse complexity, and coreference difficulty, it is hard to isolate the effect of character count per se.
\citet{baruah2025chatter} built the largest character-attribution dataset (88K character-trope pairs from 660~movies) and found that full-script context actually \emph{decreased} performance for most models relative to priors alone---evidence that models struggle to aggregate distributed character information across long contexts.
Other recent work has tested character-location tracking in fairy tales and novels~\citep{locations2025}, plot-hole detection requiring character-state consistency~\citep{flawed2025}, and character-trait coherence in generated narratives~\citep{mary2025}.
These studies consistently find that models struggle with character tracking in naturalistic settings, but none isolates the number of characters as an independent variable.

\para{Theory of mind and mental state tracking.}
\citet{xu2024opentom} introduced \opentom, a theory-of-mind benchmark with 696~narratives and 16K~questions testing both physical and psychological state tracking for multiple characters.
\citet{evolvtrip2025} proposed augmenting LLMs with external temporal knowledge graphs to compensate for failures in tracking evolving character mental states.
These works focus on the \emph{complexity} of tracked states (beliefs, intentions) rather than the \emph{number} of entities tracked, complementing our focus on scaling.

\para{Mechanistic interpretability of state tracking.}
\citet{li2025state} used probing and activation patching on the permutation-composition task to identify two algorithms transformers learn for state tracking: an Associative Algorithm (AA) that performs a hierarchical parallel scan, and a Parity-Associative Algorithm (PAA) that first computes parity, then applies associative operations.
They found that code-pretrained models preferentially learn AA, which generalizes better to longer sequences.
Their work provides a mechanistic basis for interpreting our behavioral findings: the confusion errors we observe are consistent with the overlapping distributed representations characteristic of associative state tracking.

\para{Positioning our work.}
Our study bridges the gap between controlled synthetic benchmarks~\citep{kim2023entity,hsieh2024ruler} and naturalistic evaluations~\citep{hamilton2025tldm,baruah2025chatter}.
We retain the experimental control of synthetic tasks while framing the task as narrative character tracking, and we are the first to systematically vary the number of tracked entities from 2 to 64.
Our error taxonomy (confusion vs.\ forgetting) connects behavioral measurements to mechanistic hypotheses about transformer state tracking~\citep{li2025state}.
